\documentclass[border=5pt]{standalone}
\usepackage{pgfplots}
\usepackage{pgf-pie} % 饼图：AI 代码可直接使用 \pie
\pgfplotsset{compat=1.18}
\usepgfplotslibrary{fillbetween}  % 面积图 / 堆叠面积图 / \closedcycle 填充路径
% 注：error bars 是 pgfplots 内置功能（在核心 pgfplots.errorbars.code.tex），不需要单独 \usepgfplotslibrary 加载
\usepackage{amsmath}
\usepackage{amssymb}
\usepackage{xcolor}[dvipsnames,svgnames]  % 加载标准色名表（steelblue/teal/orange/coral 等），避免 AI 用预定义色名时报 Undefined color

% 支持中文
\usepackage{fontspec}
\usepackage{xeCJK}
\setCJKmainfont{SimSun}
\setmainfont{Times New Roman}

\begin{document}

\begin{tikzpicture}
\begin{axis}[
    ybar,
    width=10cm, height=6cm,
    title={各类型项目投资金额对比},
    xlabel={项目类型},
    ylabel={投资金额（万元）},
    symbolic x coords={类型A,类型B,类型C,类型D,类型E},
    xtick=data,
    nodes near coords,
    every node near coord/.append style={font=\scriptsize, fill=none, draw=none, inner sep=1pt, anchor=south},
    legend style={at=(1.03,0.5), anchor=west, draw=black, fill=white},
    enlarge x limits=0.15,
    bar width=0.4cm,
    color={blue, red, green, orange, purple}
]
\addplot coordinates {(类型A,1200) (类型B,950) (类型C,1500) (类型D,800) (类型E,1100)};
\addlegendentry{投资金额（万元）}
\node[anchor=north west, font=\scriptsize] at (axis description cs:0.0,-0.15) {数据来源：内部统计报表 2024};
\end{axis}
\end{tikzpicture}

\end{document}